% Section 2. The validity ladder. Full specification: Supplement S7 and paper_brm/LADDER_SPEC.md.

\section{The Validity Ladder}
\label{sec:ladder}

Whether a simulated sample resembles people breaks into several questions.
Each use of the sample is a separate inference from the model's answers to a population, and each can hold or fail on its own.
Reading a persona's score, treating one answer as a case, comparing groups, estimating a group's level and scoring a scale each rest on a different property of the sample.
The validity ladder names these uses, pairs each with the property it depends on and a test against human data, and orders them from narrow to broad, beginning with a gate that tests whether a score is stable enough to read.
In these terms, a researcher who prompts a language model with personas and collects its answers to a fixed instrument has a simulated sample.
The instrument may be a multi-item scale, a single survey item or an opinion question.
Each call to the model returns one draw, and each persona under one model and one prompt framing forms a cell of $k$ draws whose average is the persona mean.

The rungs differ in what they need (Figure~\ref{fig:ladder}).
Repeated draws of each persona are enough for the gate.
Levels 1 and 4 require a multi-item scale and individual human answers on it, and levels 2 and 3 require personas built on attributes the reference records and a reference that measures the resulting groups precisely.
Where a study cannot meet a rung's requirements, that rung is not run and its claims are not made.
Verdicts are given per model, because a pooled verdict can hide errors running in opposite directions.
Passing a test means it found no failure at the precision the data allow, and a test the reference is too imprecise to support is recorded as not read.
Only the gate is a prerequisite, and a failure at one level leaves the others to run.
The gate sets $k$ for levels 2 and 3.
A group comparison of a multi-item total also assumes that the scale measures the same construct in each group, and a pass at level 2 or 3 on a total is valid only beside a pass on the invariance rule of level 4.
All pass rules and thresholds serve as defaults a user may change.
Supplement S7 gives the full specification.

\begin{figure}[!ht]
\input{figures/fig1_ladder_caption}
{\centering\includegraphics[width=\linewidth]{fig1_ladder.pdf}\par}
\input{figures/fig1_ladder_note}
\end{figure}
\FloatBarrier

\subsection{Gate: Regeneration Stability}
\label{sec:gate-method}

Two calls to a model with the same persona prompt rarely return the same answer, and a study that reads one generation as a persona's answer reads the persona's score plus noise.
The gate asks how many draws must be averaged before that noise is small enough for the score to be read.
Following a generalizability study \citep{cronbach1972dependability,brennan2001gtheory}, the variation splits into $\sigma^2_p$, the spread between personas, and $\sigma^2_r$, the spread among draws of one persona.
Averaging $k$ draws gives a standard error and a dependability of
\begin{equation}
\mathrm{SE}(k) = \sqrt{\sigma^2_r / k}, \qquad \phi(k) = \frac{\sigma^2_p}{\sigma^2_p + \sigma^2_r / k}.
\label{eq:phi}
\end{equation}
Means of $k$ draws may be read once $\mathrm{SE}(k)$ falls to 0.25 reference standard deviations in each prompt framing, and the recommended $k$ takes it to 0.125.
Single draws may be read only if the draw-to-draw standard deviation is itself no larger than 0.25 reference standard deviations.
Dependability is reported beside every verdict and required ($\phi \ge .80$) only for studies that compare or rank individual personas.
Real demographic cells differ little, and a faithful simulator shows the same small differences.
Passing supports reading the $k$-draw mean as the model's answer for that persona, leaving the correctness of that answer to the levels above, and failing rules out persona-level readings until more draws are averaged.

\subsection{Level 1: Individual Coherence}
\label{sec:l1-method}

Exemplar vignettes, screener test items and synthetic respondents treat single draws as individual cases.
Real people with the same total answer in different patterns, most endorsing the common items before the rare ones and a minority answering in unusual ways, while a simulator can return each total in the same textbook pattern.
Level 1 fits a graded response model \citep{samejima1969graded} to the human reference with survey weights and gives every respondent and every draw the standardized person-fit statistic $l_z^*$ \citep{drasgow1985appropriateness,snijders2001personfit,sinharay2016personfit}, low for an unusual pattern and high for an unusually regular one.
Each draw is set against real respondents with the same total.
The \emph{misfit ratio} divides the share of draws more unusual than 95\% of matched respondents by the 5\% expected, and the \emph{overfit ratio} does the same for draws more regular than 95\% of them.
In a real population both equal 1.
A model passes in a framing when the 90\% intervals of both ratios fall between $1/\tau$ and $\tau$, where $\tau$ is the smallest of 1.5, 2 and 3 that covers every clear departure from 1 shown by real subgroups of the reference, with a floor of 1.5 set by the controls.
A model with too few patterns in both tails is \emph{compressed}.
Passing supports single draws as individual cases, and a compressed failure rules out any use resting on the spread of real presentations.

\subsection{Level 2: Subgroup Fidelity}
\label{sec:l2-method}

Many persona studies compare groups, reporting an income gradient in a symptom score or a partisan gap in an attitude.
Level 2 asks whether the gap between two groups has the population's size and direction.
The simulated gap $g$ compares personas that differ on one attribute and match on the rest, and the population gap $\gamma$ is computed the same way among people in the reference who match on the other attributes.
Level 2 reads the ratio $\rho = g/\gamma$, where 1 reproduces the population gap, 2 doubles it, 0 erases it and a negative ratio reverses it, with a Fieller interval \citep{fieller1954interval}.
The contrasts tested together on one model form its family, and each interval is widened by a Bonferroni correction over the family.

A contrast is \vd{kept} when its whole interval lies between 0.75 and 1.25, placing the simulated gap within a quarter of the real gap's size.
A contrast is \vd{not kept} when the interval lies wholly outside that range and \vd{unresolved} when it overlaps it.
A contrast that is not kept is named by its ratio, with ratios above 1.25 steepened, from 0.25 to 0.75 attenuated, near zero missing and negative reversed.
A contrast is \vd{not read} in two cases.
The first is when the reference cannot tell the contrast's population gap from zero.
The second is when the reference measures the gap too imprecisely for any simulated gap to be read kept, unless the contrast is already clearly not kept.
A model passes level 2 when at least one contrast is read and every contrast read is kept, fails when any is not kept, and is otherwise unresolved.
A pass supports reading the size and direction of group differences from the sample, leaving either group's level to level 3, and a failure rules out the group comparisons it names.

The second case depends on the reference alone and can be checked before any model is run, a check the ladder calls step 0.
A contrast is certifiable when a faithful simulator with no sampling noise, whose gap equals the population gap, would be read kept at the family's Bonferroni level (Supplement S7).
When a family has no certifiable contrast and its verdict is unresolved, it is reported as \vd{reference cannot certify}, and passes and failures stay as they are.
Run on a planned reference, step 0 lists the gaps a study could certify, and a researcher can enlarge the reference or narrow the claims before generating.

\subsection{Level 3: Population Calibration}
\label{sec:l3-method}

Level 3 serves studies that use a simulated sample in place of a survey to estimate a group's level, such as the prevalence of a condition or the share agreeing with a policy.
A persona grid holds one persona for each combination of attributes, while a real group contains those combinations in unequal shares.
Each persona cell $c$ in a group $G$ is therefore weighted by its share $p_c$ of that group in the reference, and the residual for the group is
\begin{equation}
R_G = \sum_{c \in G} p_c\,(s_c - y_c),
\end{equation}
where $s_c$ is the persona mean and $y_c$ the mean of real respondents who match that persona \citep{holtsmith1979poststrat}.
Draw error within cells and a design-based jackknife of the survey together give its interval.
A group row passes at tolerance $\delta$ when its 90\% interval lies within $\pm\delta$, by two one-sided tests \citep{schuirmann1987tost,lakens2018tutorial}, with $\delta$ named by the researcher as the smallest difference that matters for the use.
Half a reference standard deviation is the default and 0.2 standard deviations the strict alternative.
Because a correct mean can hide a wrong prevalence, the share at or above the instrument's cut point enters as a second row, and every row must pass.
A pass supports using the simulated level in place of the group's level, within $\delta$, for the prompt, persona grid and survey period tested.

\subsection{Level 4: Structural Fidelity}
\label{sec:l4-method}

Level 4 serves studies that sum the items of a scale into one total and compare groups on it.
In people the items of a validated scale rise and fall together, and the total measures one thing.
A simulator can return plausible totals whose items do not.
Level 4 treats the draws of one model in one framing as one population, resampled by persona, and fits a one-factor model to the polychoric item correlations \citep{olsson1979polychoric}.
The fit is checked against three rules, numbered as in the full specification in Supplement S7.
Rule R1, a general factor, requires every loading to be at least .30 and the first eigenvalue at least three times the second, or the reference's own value where the reference falls short.
Rule R2, loadings like people's, requires Tucker's congruence of at least .95 \citep{lorenzoseva2006congruence} and a root mean square loading difference of at most .10, each at the conservative end of its 90\% interval.
Rule R4, the same scale in every group, requires the factor to be invariant across the groups the reference records wherever it is in people, read on the RMSEA of each invariance step \citep{savalei2024rmsea} against a margin of .08.
A model passes level 4 when it passes all three rules in every framing.
R1 and R2 support scoring the total as one score with item weights like people's, and R4 adds that a group comparison of the total means what it means in people.
The specification's R3, whether the factor holds among draws of one persona, is a diagnostic and does not decide the verdict.
